
\item Let $a,b\in\mathbb{R}$ and consider the matrices

$$
A=
\begin{bmatrix}
8&-2&1&3\\
0&a^2-4&a-2&6\\
0&0&a^2-4&b\\
0&0&0&a
\end{bmatrix},
\qquad
B=
\begin{bmatrix}
b\\1\\2\\2b
\end{bmatrix}.
$$

\begin{enumerate}[label=\alph*)]

\item Let $(S)$ be the system whose augmented matrix is $[A\mid B]$.
Discuss $(S)$ depending on $a$ and $b$.

\begin{answer}
    The augmented matrix is
    $$
    \left[
    \begin{array}{cccc|c}
    8&-2&1&3&b\\
    0&a^2-4&a-2&6&1\\
    0&0&a^2-4&b&2\\
    0&0&0&a&2b
    \end{array}
    \right].
    $$

    The coefficient matrix is upper triangular. Its diagonal entries are
    $$
    8,\quad a^2-4,\quad
    a^2-4,\quad a.
    $$

    Since
    $$
    a^2-4=(a-2)(a+2),
    $$
    the exceptional values are
    $$
    \boxed{a=-2,\quad a=0,\quad a=2.}
    $$

    {\color{red}
    \textbf{Note:} A vanishing potential pivot requires a separate discussion.
    It does not automatically imply that the system is consistent and indeterminate}

    \textbf{Case 1:
    $a\notin\{-2,0,2\}$.}

    Every diagonal entry is nonzero.
    Therefore, $A$ has four pivots:
    $$
    r(A)=4.
    $$

    Since $[A\mid B]$ has four rows,
    $$
    r(A)\leq r([A\mid B])\leq4.
    $$

    Consequently,
    $$
    r(A)=r([A\mid B])=4=n.
    $$

    The system is consistent and determined for every $b$.

    \textbf{Case 2: $a=0$.}

    Substituting $a=0$ gives
    $$
    \left[
    \begin{array}{cccc|c}
    8&-2&1&3&b\\
    0&-4&-2&6&1\\
    0&0&-4&b&2\\
    0&0&0&0&2b
    \end{array}
    \right].
    $$

    The coefficient matrix has three pivots, so
    $$
    r(A)=3.
    $$

    If $b=0$, the last row becomes entirely zero:
    $$
    r(A)=r([A\mid B])=3<4.
    $$

    Therefore, the system is consistent and indeterminate, with degree
    $$
    p=4-3=1.
    $$

    If $b\neq0$, the last row represents
    $$
    0=2b,
    $$
    which is a contradiction.

    Thus,
    $$
    r(A)=3<4=r([A\mid B]).
    $$

    The system is inconsistent.

    \textbf{Case 3: $a=2$.}

    Substitution gives
    $$
    \left[
    \begin{array}{cccc|c}
    8&-2&1&3&b\\
    0&0&0&6&1\\
    0&0&0&b&2\\
    0&0&0&2&2b
    \end{array}
    \right].
    $$

    The last three equations are
    $$
    \begin{cases}
    6x_4=1,\\
    bx_4=2,\\
    2x_4=2b.
    \end{cases}
    $$

    From the first equation,
    $$
    x_4=\frac16.
    $$

    From the third equation,
    $$
    x_4=b.
    $$

    Therefore,
    $$
    b=\frac16.
    $$

    However, substituting $x_4=\frac16$ into the
    second equation gives
    $$
    \frac{b}{6}=2
    \implies b=12.
    $$

    This contradicts
    $b=\frac16$.

    Hence the system is
    inconsistent for every $b$.

    More precisely,
    $$
    r(A)=2<3=r([A\mid B]).
    $$

    \textbf{Case 4: $a=-2$.}

    Substitution gives
    $$
    \left[
    \begin{array}{cccc|c}
    8&-2&1&3&b\\
    0&0&-4&6&1\\
    0&0&0&b&2\\
    0&0&0&-2&2b
    \end{array}
    \right].
    $$

    The last two equations are
    $$
    \begin{cases}
    bx_4=2,\\
    -2x_4=2b.
    \end{cases}
    $$

    From the second equation,
    $$
    x_4=-b.
    $$

    Substitution into the first gives
    $$
    -b^2=2.
    $$

    Equivalently,
    $$
    b^2=-2.
    $$

    Since $b\in\mathbb{R}$, this is impossible.

    Thus the system is inconsistent for every real value of $b$.

    In this case,
    $$
    r(A)=3<4=r([A\mid B]).
    $$

    \textbf{Conclusion}

    $$
    \boxed{
    \begin{array}{ll}
    a\notin\{-2,0,2\}:&
    \text{SCD, for every }b,\\
    a=0,\ b=0:&
    \text{SCI, degree }1,\\
    a=0,\ b\neq0:&
    \text{SI},\\
    a\in\{-2,2\}:&
    \text{SI, for every }b.
    \end{array}
    }
    $$

    {\color{red}
    \textbf{Important:} To discuss a system, we must determine the ranks of the coefficient
    and augmented matrices. Computing the general solution is not required.}
\end{answer}

\item Determine the values of $a$
and $b$ for which the matrix
$A-5I_4$ is invertible.

\begin{answer}
    Subtracting $5I_4$ changes only the diagonal entries:
    $$
    A-5I_4=
    \begin{bmatrix}
    3&-2&1&3\\
    0&a^2-9&a-2&6\\
    0&0&a^2-9&b\\
    0&0&0&a-5
    \end{bmatrix}.
    $$

    This matrix is upper triangular. Its diagonal entries are
    $$
    3,\quad a^2-9,\quad
    a^2-9,\quad a-5.
    $$

    Since it is a square matrix of order $4$,
    $$
    A-5I_4\text{ invertible}
    \iff r(A-5I_4)=4.
    $$

    All four diagonal entries must therefore be nonzero.

    Since
    $$
    a^2-9=(a-3)(a+3),
    $$
    the required conditions are
    $$
    a\neq-3,\qquad
    a\neq3,\qquad
    a\neq5.
    $$

    Hence,
    $$
    \boxed{
    a\in\mathbb{R}
    \setminus\{-3,3,5\},
    \qquad b\in\mathbb{R}.
    }
    $$

\end{answer}

\item For the values determine in part b), determine, with
justification, whether there exists a nonzero column matrix
$C\in M_{4\times1}(\mathbb{R})$ such that $AC=5C$.

\begin{answer}
    We seek a nonzero column matrix $C$ satisfying
    $$
    AC=5C.
    $$

    Since
    $$
    5C=(5I_4)C,
    $$
    we obtain
    $$
    \begin{aligned}
    AC&=5I_4C\\
    AC-5I_4C&=0\\
    (A-5I_4)C&=0.
    \end{aligned}
    $$

    This is a homogeneous linear system whose
    coefficient matrix is $A-5I_4$.

    For the values found in part b), this matrix is invertible.

    Multiplying on the left by its inverse gives
    $$
    \begin{aligned}
    (A-5I_4)C&=0\\
    (A-5I_4)^{-1}
    (A-5I_4)C&=0\\
    I_4C&=0\\
    C&=0.
    \end{aligned}
    $$

    Therefore, the only solution is the zero column.

    Consequently,
    $$
    \boxed{
    \text{No nonzero such column exists.}
    }
    $$

    {\color{red}
    \textbf{Note:} The expression
    $AC=5C$ can be rewritten
    as $(A-5I_4)C=0$.

    An invertible coefficient matrix gives a homogeneous
    system with a unique solution, necessarily the zero vector.
    }
\end{answer}

\end{enumerate}
